\section{Correlations, dissipation, and rational feasibility}
\label{sec:a-corr-energy}

The product structure in \cref{thm:a-design-independent} permits arbitrary
correlations within a cycle, but requires independence between cycles.
Global resistance correlations separate three tasks: individual arc
validation remains linear programming on a cactus
(\cref{thm:a-design-global-capacity}); maximizing total edge flow or
minimizing dissipation becomes strongly hard; and maximizing dissipation
admits an additive convex algorithm on every graph. Exact physical
constraints introduce a further distinction between rational and algebraic
parameter witnesses. Unless stated otherwise, this section uses the
symmetric quadratic law and fixed rational balanced nominations.

\subsection{A shared parameter and a curved attainable region}

\begin{example}[Two correlated cycles]
\label{ex:a-corr-shared}
Join two source-to-sink triangles by a bridge of resistance one. Supply one
unit at the first entrance and withdraw one at the last exit. The first
triangle has a two-edge path with resistances $1/2,1/2$; the second has a
two-edge path with resistances $2,2$. Both direct edges have the same
resistance $\theta\in[1,4]$. This is a simple degree-three cactus, and its
uncertainty set is a rational polytope in the actual edge resistances.
Equal path drops and conservation give the two direct flows
\[
 x_1=\frac1{1+\sqrt\theta},\qquad
 x_2=\frac2{2+\sqrt\theta},\qquad
 x_2=\frac{2x_1}{1+x_1},\quad \frac13\le x_1\le\frac12.
\]
All edge flows are positive. The last graph is strictly concave and
nonaffine, so the projected attainable region is nonconvex; hence so is
the full flow region. More explicitly, the midpoint of its endpoints is
$(5/12,7/12)$, whereas the graph at $5/12$ has ordinate $10/17$, a
difference of $1/204$.

The linear objective $F=x_2-x_1$ has value $1/6$ at both parameter
endpoints. Writing $s=\sqrt\theta$ gives
$F=s/((s+1)(s+2))$ and
\[
 \frac{dF}{ds}=\frac{2-s^2}{(s^2+3s+2)^2}.
\]
Its unique maximum is therefore at $\theta=2$, where
$F=3-2\sqrt2$. The gain over either endpoint is
\[
 \frac{17}{6}-2\sqrt2
 =\frac{1}{36(17/6+2\sqrt2)}>\frac1{216};
\]
the last inequality follows from $\sqrt2<19/12$. Measuring one edge of
the first alternate path and the second direct edge gives the positive
unit-coefficient objective $1+F$, with the same interior optimum.
Multiplying all resistances by two makes the fixed resistances integral
without changing any flow. This example disproves a shared-parameter
endpoint rule; by itself it makes no complexity-hardness assertion.
\end{example}

\subsection{Bounded-data hardness for total flow and minimum dissipation}

For a physical state define its constitutive dissipation by
\begin{equation}
 D(\beta)=\sum_{e\in E}\beta_e|x_e(\beta)|^3
         =b^\top\pi(\beta).
 \label{eq:a-corr-dissipation}
\end{equation}
The equality follows by multiplying $A^\top\pi=g(x)$ by $x$ and using
$Ax=b$. Under unit source/sink nominations it is the prescribed terminal
potential difference. If potentials represent squared gas pressures,
this is a squared-pressure difference; the identity does not identify
$D$ with literal compressor power.

\begin{theorem}[Globally correlated cactus hardness]
\label{thm:a-corr-hardness}
On connected simple cacti of maximum degree three, with unit source/sink
nominations and a global rational resistance polytope, each of the
following problems is strongly $\NP$-hard: maximizing
$T=\sum_e x_e$, and minimizing $D$. This holds with resistances in $[1,3]$,
bounded integer coefficients in the polytope description, an acyclic
orientation, and strictly positive flows on every edge. Thus $T$ also
sums all absolute edge flows, with coefficient one on every edge.
Absolute-error $1/256$ value approximation is $\NP$-hard for either
problem, as is returning a rational admissible profile within $1/256$
of the appropriate optimum.

On the explicit families constructed below, the weak existential tests
$T\ge\tau$ and $D\le\tau$ are $\NP$-complete, and their strict universal
complements are $\coNP$-complete. The universal test $T\le\tau$ is also
$\coNP$-complete on the total-flow family. No membership statement is
asserted for arbitrary correlated cactus objectives.
\end{theorem}
\begin{proof}
Use an unweighted Max-Cut instance $H$ with $n_H$ vertices, $m_H\ge1$
edges, and integer target $1\le K\le m_H$. The standard binary cut
encoding is classical; see, for example,
\citet[Section 1.1]{delpia2017-miqp-np}. Trivial source cases can be
replaced by a one-edge comparison graph with target one for a yes case,
or a comparison triangle with target three for a no case.
Introduce $\theta_i\in[0,1]$. For every comparison edge $ij$, use two
unit-through-flow triangles. Each triangle has a two-edge path and an
alternate direct edge of resistance two. In the plus triangle, both path
edges have resistance $2+\theta_i-\theta_j$; in the minus triangle, both
have resistance $2-\theta_i+\theta_j$.

If the resistance on each path edge is $r$, its path flow $f(r)$ obeys
$2rf(r)^2=2(1-f(r))^2$. The positive solution and the triangle's terminal
drop are
\begin{equation}
 f(r)=\frac1{1+\sqrt r},\qquad
 e(r)=\frac{2r}{(1+\sqrt r)^2}.
 \label{eq:a-corr-triangle}
\end{equation}
The three edge flows sum to $1+f(r)$. Their dissipation is $e(r)$:
the common path drop is multiplied by total through-flow one. Direct
differentiation gives
\[
 f''(r)=\frac{1+3\sqrt r}{4r^{3/2}(1+\sqrt r)^3}>0,
 \qquad e'(r)=\frac2{(1+\sqrt r)^3}.
\]
\[
 e''(r)=-\frac3{\sqrt r(1+\sqrt r)^4}<0.
\]
Consequently $f(2+d)+f(2-d)$ is even and convex on $[-1,1]$, whereas
$e(2+d)+e(2-d)$ is even and concave there.

For total flow, join all $2m_H$ triangles in a chain with $2m_H-1$
bridges of resistance two. Place a path of $n_H$ additional bridges
before the first triangle, with actual resistance coordinates
$r_i\in[1,2]$, and set $\theta_i=r_i-1$. Impose the equalities
\[
 \beta_{ij,+,1}=\beta_{ij,+,2}=2+r_i-r_j,
 \qquad
 \beta_{ij,-,1}=\beta_{ij,-,2}=2-r_i+r_j.
\]
All other resistances are fixed at two. These equations and interval
bounds describe an injective affine image of the cube in the actual
physical resistance vector; no projection encoding is assumed. Their
coefficients and constants belong to a fixed integer alphabet.
Orient every route from the first source to the final sink and put
nominations $+1,-1$ there. Every bridge carries one, each triangle
carries one in total, and \cref{eq:a-corr-triangle} makes every edge
flow strictly positive. The graph is simple, is a cactus of maximum
degree three, and has an acyclic orientation. Its vertex and edge
counts are $6m_H+n_H$ and $8m_H-1+n_H$, so its cycle rank is $2m_H$.

At a binary cube vertex, one triangle pair contributes
\[
 A=2f(2)=2\sqrt2-2 \quad\hbox{if uncut},\qquad
 B=f(1)+f(3)=\frac{\sqrt3}{2} \quad\hbox{if cut}.
\]
Put $\kappa=B-A$ and $C=4m_H-1+n_H+m_HA$. The full objective is
\[
 T(\theta)=4m_H-1+n_H+
   \sum_{ij\in E(H)}[f(2+\theta_i-\theta_j)+f(2-\theta_i+\theta_j)].
\]
It is convex on the cube. A convex function on a cube has a maximizing
vertex: with all other coordinates fixed, at least one endpoint of an
interval has value at least its interior value; apply this successively.
At binary vertices $T=C+\kappa\operatorname{cut}(\theta)$, whence
\begin{equation}
 \max T=C+\kappa\operatorname{MaxCut}(H).
 \label{eq:a-corr-cut-flow}
\end{equation}
The rational square comparisons
$\sqrt3>173/100$ and $\sqrt2<283/200$ give
$\kappa>7/200>1/32$.

For minimum dissipation replace the prefix by $2n_H$ bridges, in pairs
$r_i=1+\theta_i$ and $s_i=2-\theta_i$. Impose $r_i+s_i=3$, and use the
same triangle correlations in terms of $r_i$. This is again an
injective affine cube in actual resistances, all in $[1,3]$. The graph
now has $6m_H+2n_H$ vertices and $8m_H-1+2n_H$ edges; its rank, degree,
simplicity, orientation, and positivity properties are unchanged. Each
prefix pair contributes constant dissipation three. This cancellation
is necessary because a single unit-flow bridge would contribute its
variable resistance to the objective.

At binary parameters the pair energies are
\[
 A_E=2e(2)=24-16\sqrt2,
 \qquad B_E=e(1)+e(3)=\frac{13}{2}-3\sqrt3.
\]
Put $\kappa_E=A_E-B_E$ and $C_E=3n_H+4m_H-2+m_HA_E$.
The total energy is a sum of the concave pair responses and constant
bridge energies, so it has a minimizing cube vertex by the reverse
endpoint argument. Therefore
\begin{equation}
 \min D=C_E-\kappa_E\operatorname{MaxCut}(H).
 \label{eq:a-corr-cut-energy}
\end{equation}
The same rational square comparisons give
$\kappa_E=35/2-16\sqrt2+3\sqrt3>1/20>1/32$.

We give the threshold encoding explicitly. For total flow choose rational
$\tau_T$ within $1/512$ of
\[
 Q_T=C+\kappa(K-1/2)
 =4m_H-1+n_H+(m_H-K+1/2)A+(K-1/2)B.
\]
Dyadic enclosures of $\sqrt2,\sqrt3$ of width at most
$1/[8192(m_H+1)]$ suffice: substitute any enclosed rational values;
the sum of the absolute sensitivity coefficients is at most
$2m_H+m_H/2$. For energy choose rational $\tau_E$ within $1/512$ of
\[
 Q_E=C_E-\kappa_E(K-1/2)
 =3n_H+4m_H-2+(m_H-K+1/2)A_E+(K-1/2)B_E.
\]
Width $1/[65536(m_H+1)]$ suffices, since the corresponding coefficient
sum is at most $19m_H$. Dyadic bisection needs $O(\log(m_H+1))$ bits;
taking the first power of two satisfying either bound makes the common
denominator $O(m_H+1)$. Both thresholds then have polynomially bounded
numerators in $m_H+n_H$ and polynomially bounded denominators. Thus all
numerical input has polynomial unary encoding length, including the
thresholds, establishing strong hardness.

In a yes case, \cref{eq:a-corr-cut-flow} exceeds $\tau_T+7/512$ and
\cref{eq:a-corr-cut-energy} is below $\tau_E-7/512$. In a no case the
respective optima are below $\tau_T-7/512$ and above $\tau_E+7/512$.
Indeed half of either cut reward exceeds $1/64=8/512$, and the rational
threshold incurs at most $1/512$ error. These strict gaps prove all
hardness statements for the weak tests and their complementary universal
tests, including robust $T\le\tau_T$.

For membership on these families, an extremum occurs at binary
$\theta$, hence at integer resistances. Its value lies in the fixed
field $\Q(\sqrt2,\sqrt3)$ and is exactly comparable with any input
rational threshold in polynomial time. A binary cut is thus a polynomial
certificate for existential attainment, or for strict violation of a
universal bound. This also proves the claimed universal memberships.
It does not supply such a certificate for arbitrary sums of unrelated
algebraic cycle responses.

A value estimate within $1/256=2/512$ distinguishes the two cases.
Suppose instead that an algorithm returns a rational admissible profile
within that loss. In a yes total-flow case its value exceeds
$\tau_T+5/512$, whereas every no-case profile is below
$\tau_T-7/512$. For minimum energy the signs reverse. Evaluate the
returned profile by enclosing the individual square roots in
\cref{eq:a-corr-triangle}, adding the rational bridge contributions,
and refining to total error at most $1/512$. All arguments of the roots
belong to $[1,3]$; the formulas have uniformly bounded derivatives there.
Thus $O(\log(m_H+1))$ additional accuracy bits beyond the output encoding
suffice. A polynomial-time output algorithm has polynomial output
length, so this evaluation is polynomial-time and separates the cases
without exact comparison of a general radical sum.
\end{proof}

The transported amount in this construction is fixed at one. Total edge
flow counts it again on every traversed edge; it is not end-to-end
throughput. The cycle rank grows, and neither a fixed-rank lower bound nor
constant relative-error or normalized-objective hardness follows.
The reduction uses the classical convex/concave cube-vertex mechanism
and Max-Cut. Its specific content is the bounded passive-network
realization under global resistance correlations. For comparison,
\citet[Corollary 4 and Theorem 11]{klimm2026-approximating-the-network-design-problem}
study conductance investment with installation costs, obtain convexity
without fixed costs, and give a series-parallel approximation scheme when
all variable conductance costs are positive. Linear correlations in
resistance coordinates do not become linear correlations in conductance
coordinates; those conclusions concern different feasible sets and
objectives.

\subsection{Maximum dissipation on arbitrary graphs}

The opposite energy direction has a classical variational explanation.
Let
\begin{equation}
 V(\beta)=\min_{Ax=b}\sum_e\frac{\beta_e|x_e|^3}{3},
 \qquad D(\beta)=3V(\beta).
 \label{eq:a-corr-primitive}
\end{equation}
The second equality is \cref{eq:a-corr-dissipation} and
\cref{thm:a-pre-existence}. As an infimum of affine functions of
$\beta$, $V$ is concave. For linear two-terminal networks, concavity
in all edge resistances was proved by
\citet{shannon1956-concavity}; the present conclusion follows from the
same variational structure with a cubic primitive. We give an explicit conic formulation and its
bit-complexity and exact-feasibility details. The energy principle and its
homogeneous dual are
classical (see \cref{sec:a-pre-preliminaries} and
\citet[Section 3, equation (7), and Lemma 3.4, equation (8)]{gro2017-algorithmic-results-for-potential-based});
perspectives and the hyperbolic-cone construction are also established
\citep[Section 3.2.6]{boyd2004-convex}
\citep[Section 2.3]{lobo1998-socp}.

\begin{theorem}[Additive maximum dissipation and rational profiles]
\label{thm:a-corr-energy-max}
Let $G$ be any connected graph with fixed rational balanced nominations.
Let $P$ be a nonempty bounded rational polytope of resistances, given by
rational inequalities and explicit positive bounds
$\beta_L\le\beta_e\le\beta_U$. Maximum physical dissipation over $P$
has an exact formulation with two second-order cones per edge.
For every rational $\eps>0$, a rational $\widehat\beta\in P$ satisfying
\[
 D(\widehat\beta)\ge\max_{\beta\in P}D(\beta)-\eps
\]
is computable in time polynomial in input length and
$\max\{0,\log(1/\eps)\}$. The theorem imposes no additional physical
capacity or potential filters.
\end{theorem}
\begin{proof}
For $\beta>0$, the scalar maximum of
$dx-\beta|x|^3/3$ occurs at
$x=\sgn(d)\sqrt{|d|/\beta}$ and conjugate
$2|d|^{3/2}/(3\sqrt\beta)$, also at $d=0$. Fenchel's inequality gives
\begin{equation}
 V(\beta)=\max_{\pi:\pi_{v_0}=0}
 \left[b^\top\pi-\frac23\sum_e
     \sqrt{\frac{|(A^\top\pi)_e|^3}{\beta_e}}\right].
 \label{eq:a-corr-dual}
\end{equation}
For completeness, weak duality follows by summing Fenchel's inequality
at any conserved flow. The physical potentials from
\cref{thm:a-pre-existence} give equality on every edge, proving attained
strong duality. Normalization changes neither edge drops nor the balanced
pairing $b^\top\pi$. The conjugate is the convex perspective
$\beta\, (2/3)|d/\beta|^{3/2}$, so the bracket is jointly concave in
$(\pi,\beta)$. Combining its maximum over $\pi$ with a maximum over
$\beta\in P$ interchanges only two maxima.

For each edge introduce $z_e,u_e,t_e\ge0$, put
$d_e=(A^\top\pi)_e$, and impose
\begin{equation}
 z_e\ge d_e,\quad z_e\ge-d_e,\quad
 z_e^2\le t_eu_e,\quad u_e^2\le\beta_e z_e.
 \label{eq:a-corr-cones}
\end{equation}
These are two rotated second-order cones and linear inequalities;
for example $z^2\le tu$ with $t,u\ge0$ is
$\|(2z,t-u)\|_2\le t+u$, with rational coefficients.
If $z>0$, multiplying gives $z^4\le\beta zt^2$, hence
$z^3\le\beta t^2$. Conversely this last inequality with $z>0$
forces $t>0$, and $u=z^2/t$ satisfies both cones. At $z=0$ choose
$u=0$ and any $t\ge0$; conversely the second cone forces $u=0$.
If $t=0$, the first cone forces $z=0$. Thus eliminating $z,u$ is
exactly $t\ge\sqrt{|d|^3/\beta}$, including every boundary case.
Maximize
\begin{equation}
 J=b^\top\pi-\frac23\sum_e t_e
 \quad\hbox{subject to }\beta\in P,\quad\pi_{v_0}=0,
 \quad\cref{eq:a-corr-cones}.
 \label{eq:a-corr-socp}
\end{equation}
Its optimum is $\max_P V$, whose triple is the requested dissipation.

We next justify polynomial-bit exact rational feasible output, rather
than assuming a general exact SOCP solver. If $b=0$, every physical flow
and dissipation is zero, and rational LP returns any profile in $P$.
The connected edgeless case is included. Otherwise put
$M=\tfrac12\sum_v|b_v|>0$ and $m=|E|\ge1$. Orient each nonzero physical
edge downhill in potential. This orientation is acyclic and its
nonnegative flow decomposes into source-to-sink paths of total weight
$M$; one may add a supersource and supersink and repeatedly remove such
a path. Hence $|x_e|\le M$. A simple path from $v_0$ yields
\[
 |\pi_v|\le m\beta_U M^2.
\]
At a physical dual optimum, cone variables may be chosen as
\begin{equation}
 z_e=\beta_e|x_e|^2,\qquad u_e=\beta_e|x_e|,
 \qquad t_e=\beta_e|x_e|^3.
 \label{eq:a-corr-physical-lift}
\end{equation}
Both cones are tight and these coordinates are uniformly bounded.
Moreover, evaluating the primitive energy at the other profile's
minimizer gives
$|V(\beta)-V(\beta')|\le(M^3/3)\|\beta-\beta'\|_1$.
Thus $V$ attains its maximum on compact $P$.

Here are explicit compact bounds that preserve a global optimizer. Set
$r=\min\{1,\beta_L\}/2$ and use
\[
 \begin{split}
 |\pi_v|&\le m\beta_U M^2+1,\qquad
 0\le z_e\le\beta_U M^2+2,\\
 0\le u_e&\le\beta_U M+r+1,\qquad
 0\le t_e\le\beta_U M^3+2/r+1.
 \end{split}
\]
Every bound has polynomial rational encoding length and contains
\cref{eq:a-corr-physical-lift}. At a relative-interior point $\beta^0$
of $P$, use $\pi=0,z_e=1,u_e=r,t_e=2/r$. Then
$z_e>|d_e|$, $t_eu_e-z_e^2=1$, and
$\beta_e z_e-u_e^2\ge\beta_L-r^2>0$; all added upper and lower bounds
are strict as well.

We construct $\beta^0$ and a quantitative interior ball. For each row
$R_i\beta\le c_i$ of $P$, maximize its slack by rational LP. A zero
maximum identifies a universally tight row. For each other row retain
a polynomial-bit rational maximizing vertex with positive slack.
Average these witnesses (or take any feasible rational point if there
are no such rows). All nonuniversal rows are then strict at the average.
Basic-solution determinant bounds give polynomial encoding lengths for
each vertex, their average, and every positive rational slack. The
universally tight rows define the affine hull: within their solution
space the strict nonuniversal rows leave a neighborhood of the average,
so no further affine equality can hold throughout $P$.

Rational Gaussian elimination expresses this hull by
$\beta=\beta^0+Q\xi$, choosing free coordinates among original resistance
coordinates. The entries of $Q$ and the chosen center coordinates have
polynomial bit length; the free coordinates have bounds inherited from
$P$. Remove the fixed potential coordinate. Substitution leaves a
compact convex region $K$ in the remaining coordinates, with every
nonconstant constraint a rational linear or quadratic inequality
$p_j(y)\ge0$ strict at a rational center $y^0$. Universal equalities
and redundant constant inequalities are removed. If $P$ is a singleton,
there are simply no free resistance coordinates.

On $\|y-y^0\|_\infty\le1$, bound $\|\nabla p_j(y)\|_1$ by a rational
$G_j\ge1$ using the absolute coefficients and the coordinate magnitudes
$|y_i^0|+1$. Since degrees are at most two and there are polynomially
many coefficients, these bounds have polynomial bit length. The rational
radius
\[
 \rho=\min\left\{1,\min_j\frac{p_j(y^0)}{2G_j}\right\}>0
\]
puts the Euclidean ball $B_2(y^0,\rho)$ inside $K$, by the mean-value
inequality. The coordinate bounds give a rational outer radius about
$y^0$, for instance the sum of bounds on the absolute coordinate
differences. Both radii and the center have polynomial encoding length.
Convexity here follows from the nonnegative rotated-cone constraints,
not from an assumption that arbitrary quadratic inequalities are convex.

Membership of a rational point in $K$ is decided exactly by rational
linear and quadratic comparisons, supplying a weak membership oracle.
The negative objective $-J$ is linear on the whole reduced coordinate
space, hence is globally convex and Lipschitz with a rational bound
obtained by summing the absolute coefficients; its rational value oracle
is exact. Apply \citet[Theorem 2.5.9, p.~48]{dadush2012-integer-programming}
to this centered convex body with error $\eps/3$. This theorem returns a
rational point \emph{in} $K$, with an additive objective guarantee, in
polynomial time in these encodings and accuracy bits. Its feasibility
conclusion is essential here. Map its coordinates back to obtain rational
$\widehat\beta\in P$ exactly. Every feasible conic point satisfies
$J\le V(\widehat\beta)$ by \cref{eq:a-corr-dual}; therefore
\[
 D(\widehat\beta)=3V(\widehat\beta)
 \ge3J\ge\max_{\beta\in P}D(\beta)-\eps.
\]
The auxiliary potentials need not be physical. Neither rounding a
physical flow nor exactly evaluating a radical sum is required.
\end{proof}

\begin{corollary}[Exact cactus capacities]
\label{cor:a-corr-energy-capacity}
On a cactus, \cref{thm:a-corr-energy-max} remains valid with rational
signed arc capacities: emptiness is decidable in polynomial time, and a
nonempty instance admits the same additive algorithm with a rational
profile satisfying all original physical capacities exactly.
\end{corollary}
\begin{proof}
Apply \cref{thm:a-design-global-capacity} to obtain the exactly equivalent
rational polytope $P_{\rm cap}\subseteq P$, rejecting fixed bridge
violations or emptiness. This is the cycle interval linearization of
\citet[Proposition 4.9, Lemma 4.10, Proposition 4.11]{amann2018-deciding-robust-feasibility-and-infeasibility},
composed across blocks in \cref{thm:a-design-global-capacity}. Apply
\cref{thm:a-corr-energy-max} to $P_{\rm cap}$. Its affine-hull reduction
handles capacity equalities and singletons. Membership of the returned
profile in $P_{\rm cap}$ guarantees its actual physical capacities;
the conic auxiliary variables are used only for the objective guarantee.
\end{proof}

An exact conic representation is not an exact polynomial-time scalar
comparison algorithm. For fixed resistances and unit terminal nominations
on a cactus, $D$ is the terminal potential difference. The two reductions
in \cref{thm:a-cac-srs} therefore apply to its exact threshold comparison,
already for a singleton parameter polytope. The theorem above returns
rational near-optimal resistance parameters; it promises neither a
rational physical state nor an exact rational optimal parameter vector.

\subsection{An exact certificate for the original global design problem}

\begin{proposition}[Rational global design certificate]
\label{prop:a-corr-certificate}
Write the complete positive bounded resistance polytope as
$P=\{\beta:R\beta\le c\}$. Let rational $y,\lambda,\beta^0,\pi,t$ satisfy
\[
 \begin{gathered}
 Ay=b,\qquad \lambda\ge0,\qquad
 R^\top\lambda=w,\quad w_e=|y_e|^3/3,\\
 R\beta^0\le c,\qquad t_e\ge0,\qquad
 \beta_e^0t_e^2\ge|(A^\top\pi)_e|^3.
 \end{gathered}
\]
Then with $U=c^\top\lambda$ and
$L=b^\top\pi-(2/3)\sum_e t_e$,
\begin{equation}
 \max_{\beta\in P}D(\beta)-D(\beta^0)\le3(U-L).
 \label{eq:a-corr-certificate-gap}
\end{equation}
The certificate allows any additive potential gauge because $b$ is balanced.
All hypotheses and the gap are exactly checkable by rational arithmetic
in time polynomial in their encoded lengths.
\end{proposition}
\begin{proof}
For every $\beta\in P$, the conserved trial flow yields
$V(\beta)\le\beta^\top w\le c^\top\lambda=U$ by LP weak duality.
The conjugate inequalities and \cref{eq:a-corr-dual} yield
$L\le V(\beta^0)$. Subtract and multiply by three. The checks use only
rational equalities, inequalities, absolute cubes, and products.
\end{proof}

For an exact zero-gap example, orient a triangle with two path edges and
one direct edge from source to sink, and impose unit through-nominations.
Let $P=\{\beta\in[1,4]^3:\sum_e\beta_e=6\}$. Take
\[
 \begin{gathered}
 y=(1/2,1/2,1/2),\qquad \beta^0=(3/2,3/2,3),\\
 (\pi_s,\pi_w,\pi_t)=(3/4,3/8,0),\qquad
 t=(3/16,3/16,3/8).
 \end{gathered}
\]
Each $w_e=1/24$. Represent the sum equality by two inequalities and give
its upper direction multiplier $1/24$, with all other multipliers zero.
Then $U=6/24=1/4$. The edge drops are $(3/8,3/8,3/4)$; each conjugate
inequality is tight, and $L=3/4-(2/3)(3/4)=1/4$. This proves the global
maximum $D=3/4$ and exact optimality of the displayed original profile.
For the capacitated corollary, use the full rows of $P_{\rm cap}$ in the
certificate. The trial flow $y$ need only conserve nominations: it bounds
the unconstrained physical primitive energy of every admissible profile,
so no capacity condition on the trial flow is needed.

The certificate is sufficient for its stated gap. An arbitrary conserved
trial flow need not yield a small gap, and the proposition supplies no
general numerical producer. The accompanying exact rational verifier
checks the encoded graph, nominations, resistance bounds, all polytope
rows, and the displayed conditions. The polynomial verification bound
counts the bit lengths of expanded rational numerators and denominators.
Additional exact-string forms accepted by the \texttt{Fraction} parser,
including exponent notation, are subject to this bound after expansion,
not in their raw JSON length. Its polytope input contains no
independently interpreted operating constraints; deriving such rows from
physical capacities is a separate equivalence check. A numerical conic
producer for selected budget-constrained instances supplies candidate
certificates, whose rational verification, rather than solver status,
establishes their guarantees. This implementation is not a general
implementation of the convex-body algorithm in
\cref{thm:a-corr-energy-max}.

\subsection{Irrational feasible parameters beyond cacti}

\begin{proposition}[Capacities forcing irrational resistance]
\label{prop:a-corr-irrational}
There is a feasible rational-input quadratic design instance on a simple
series-parallel graph of global cycle rank two, with one resistance
interval and only nonnegative positive-width upper-capacity intervals,
for which every feasible resistance profile is irrational.
\end{proposition}
\begin{proof}
Use vertices $u,v,r,w,z$, a direct arc $a=(u,v)$, a connector $(u,r)$,
and the two routes $r,w,v$ and $r,z,v$, all oriented towards $v$. Set
\[
 \begin{gathered}
 \beta_a=\theta\in[4/3,11/8],\qquad \beta_{ur}=1,\\
 \beta_{rw}=\beta_{wv}=1/2,\qquad
 \beta_{rz}=\beta_{zv}=1.
 \end{gathered}
\]
Supply two at $u$ and withdraw two at $v$. Impose capacities $[0,1]$
on $a$ and $(u,r)$, and $[0,2]$ on the four remaining arcs. This simple
graph is a parallel composition of $a$ and a connector in series with
two parallel paths, hence is series-parallel. Its rank is $6-5+1=2$.
Conservation forces $x_a=x_{ur}=1$. The two $r$--$v$ path flows
$q_1,q_2$ have sum one and the same sign because their drops agree and
their resistance sums are positive. Thus both are positive. Their
resistance sums are one and two, so
\[
 q_1^2=2q_2^2=d,\qquad
 q_1=2-\sqrt2,\quad q_2=\sqrt2-1,\quad d=6-4\sqrt2.
\]
The connector adds drop one; the unit-flow direct arc therefore forces
\[
 \theta=7-4\sqrt2,\qquad \theta^2-14\theta+17=0.
\]
This is strictly between $4/3$ and $11/8$: the equivalent square-root
bounds are $\sqrt2<17/12$ and $\sqrt2>45/32$, verified by squaring.
Set $\pi_v=0$, $\pi_r=d$, $\pi_u=1+d$, and
$\pi_w=\pi_z=d/2$. These potentials and the displayed flows satisfy
every original equation and capacity. Conversely the cut saturation
forces all preceding values. The polynomial's discriminant is $128$,
which is not a rational square, so there is no rational feasible profile.
\end{proof}

A smaller equality-constrained version identifies $r$ with $u$ and
removes the connector. Set the direct interval to $[1/3,3/8]$ and require
$x_a=1$. With the same nomination and branch resistances, the forced
value is $\theta=6-4\sqrt2$, satisfying $\theta^2-12\theta+4=0$.
The bounds are again the rational square comparisons just used, shifted
by one. Potentials $\pi_u=\theta$, $\pi_v=0$, and
$\pi_w=\pi_z=\theta/2$ verify feasibility on this four-vertex, five-edge
graph. Replacing its interval by the two-point set $\{1/3,3/8\}$ makes
$x_a=1$ infeasible, while its interval hull remains feasible. Thus
operating-constraint filtering can destroy an unfiltered hull identity.
Discrete filtered feasibility is already hard on one cycle by
\cref{thm:a-bound-realization}; this smaller example illustrates the
failure directly.

By \cref{thm:a-design-global-capacity}, a nonempty rationally described
capacity-filtered cactus domain has a rational original parameter
witness. Every connected graph of global cycle rank at most one is a
cactus, making rank two the first possible rank for the preceding
obstruction. It is an output-representation obstruction, not an
$\NP$-hardness proof, and permits algebraic feasible output. Both cut
capacities in \cref{prop:a-corr-irrational} must be saturated; strict
margins have a different conclusion below. General irrational
arithmetic in potential-flow design is already discussed by
\citet[Remark 8]{klimm2026-approximating-the-network-design-problem}; the
point here is unavoidable irrationality of a feasible \emph{parameter},
not merely an irrational physical state at rational input.

\subsection{Uniform coefficient sensitivity, including zero flows}

\begin{theorem}[Linear flow perturbation bound]
\label{thm:a-corr-lipschitz}
Let a connected graph have $m\ge1$ edges and fixed balanced nonzero
nominations, and put $M=\tfrac12\sum_v|b_v|$. Consider two quadratic
coefficient profiles whose coefficients lie in
$[\beta_L,\beta_U]$, with $\beta_L>0$. Either use one symmetric
coefficient per edge, or two independent directional coefficients
$\beta_e^+,\beta_e^-$ as in \cref{subsec:a-pre-model}. If every
corresponding coefficient changes by at most $\delta$, the physical
flows satisfy
\begin{equation}
 \|x'-x\|_\infty\le\frac{2mM\delta}{\beta_L}.
 \label{eq:a-corr-linear-flow}
\end{equation}
The bound holds when flows vanish or reverse sign. Its constant is not
asserted to be optimal.
\end{theorem}
\begin{proof}
Both flows have magnitude at most $M$, by the downhill path decomposition
used in \cref{thm:a-corr-energy-max}; this also applies to positive
directional coefficients. Their difference $h=x'-x$ is a circulation.
Reorient every edge with negative $h_e$ so that $h_e\ge0$. Reorientation
swaps the positive and negative coefficients of each directional law,
for both old and new profiles. Thus the common bounds and the maximum
coefficient change $\delta$ are preserved.

Let $H=\|h\|_\infty>0$ and choose an arc $u\to v$ with $h_e=H$.
There is a directed path from $v$ to $u$ using only arcs with
$h_e\ge H/m$. Otherwise the vertices reachable from $v$ on such arcs
form a set containing $v$ but not $u$. The incoming cut contains flow
$H$, while every outgoing arc has flow strictly less than $H/m$.
There are at most $m$ outgoing arcs, so their total is strictly below
$H$, contradicting circulation balance. Choose a simple return path.
Together with $u\to v$, it forms a directed cycle $C$ on which every
$h_e\ge H/m$.

Write $f(x)=x|x|$. For $h>0$,
\begin{equation}
 f(x+h)-f(x)\ge\frac{h|x|}{2},
 \qquad f(x+h)-f(x)>0.
 \label{eq:a-corr-scalar}
\end{equation}
For $x\ge0$ or $x\le-h$, expansion gives even the lower bound
$h|x|$. In the crossing case $x=-th$, $0<t<1$, the difference between
the two sides of the weak inequality, divided by $h^2$, is
$2t^2-(5/2)t+1$. Its minimum is $7/32>0$.
Let $g'_e$ denote the new law. In the symmetric case, or in either
same-sign branch of a directional law,
\[
 g'_e(x+h)-g'_e(x)\ge\beta_L[f(x+h)-f(x)].
\]
In the crossing case the left side is
$\beta_e^{\prime+}(x+h)^2+\beta_e^{\prime-}x^2$, so the same inequality
holds. Also, at any fixed $x$,
$|g'_e(x)-g_e(x)|\le\delta x^2$, including at zero.

Potential differences telescope around $C$. Hence
\[
 \sum_{e\in C}[g'_e(x_e+h_e)-g'_e(x_e)]
 =-\sum_{e\in C}[g'_e(x_e)-g_e(x_e)]
 \le\delta\sum_{e\in C}|x_e|^2.
\]
If $H>2mM\delta/\beta_L$, every term with $x_e\ne0$ on the left
is at least $\beta_L h_e|x_e|/2$, strictly larger than
$M\delta|x_e|\ge\delta|x_e|^2$. Each term with $x_e=0$ is strictly
positive, while its right-hand contribution is zero. Summation gives a
contradiction. The case $H=0$ is immediate. This proof uses no positive
lower bound on differential resistance at zero flow.
\end{proof}

The elementary circulation argument gives an explicit bound without
differentiating the physical state. A regularized electrical sensitivity
argument improves its constant and retains separate directional
coefficients.

\begin{proposition}[A sharper coefficient perturbation bound]
\label{prop:a-corr-sharper}
Under the hypotheses of \cref{thm:a-corr-lipschitz}, define
$q_e=|\beta'_e-\beta_e|$ in the symmetric case and
$q_e=\max\{|\beta_e^{\prime+}-\beta_e^+|,
|\beta_e^{\prime-}-\beta_e^-|\}$ in the directional case. Then
\begin{equation}
 \|x'-x\|_\infty\le\frac{M}{2\beta_L}\sum_e q_e
 \le\frac{mM\delta}{2\beta_L}.
 \label{eq:a-corr-sharper-flow}
\end{equation}
\end{proposition}
\begin{proof}
Interpolate the coefficients affinely by a parameter $s\in[0,1]$, and
add $\rho x$ to every law, for a fixed rational $\rho>0$. The regularized
primitive is the original primitive plus $\rho\|x\|_2^2/2$. Its physical
flow $x^\rho(s)$ still has $|x_e^\rho(s)|\le M$, because its law remains
strictly increasing and positive at every positive flow.

Here is the needed differentiation and electrical calculation. Choose
a full-column-rank matrix $C$ with image $\ker A$ and write a conserved
flow as $x^0+Cz$. The stationarity equations are
$C^\top(g_s(x^0+Cz)+\rho(x^0+Cz))=0$. Directional quadratic laws are
jointly continuously differentiable in $x$ and their coefficients,
including $x=0$. The derivative in $z$ is $C^\top\mathcal R C$, where
$\mathcal R$ is diagonal with
\[
 r_e=2\beta_e^{\mathrm{active}}(s)|x_e^\rho(s)|+\rho>0.
\]
At zero either directional coefficient gives the same value $r_e=\rho$.
The matrix is positive definite, so the implicit-function theorem and
uniqueness make $x^\rho(s)$ continuously differentiable. If $\ker A=0$,
the flow is independent of all coefficients and the conclusion is
immediate. The equations may be differentiated locally beyond either
endpoint because all coefficients remain positive in a neighborhood.

At fixed $s$, let $a_e$ be the incidence column of edge $e$, and let
$j^{(e)}$ be the electrical flow with resistances $r_a$ and nomination
$a_e$. Thus $Aj^{(e)}=a_e$ and $\mathcal Rj^{(e)}=A^\top v$ for some
$v$. Existence and uniqueness follow by minimizing
$\sum_a r_a j_a^2/2$ on that affine conservation space, exactly as in
\cref{thm:a-pre-existence}. Its nonzero flow is downhill, so unit
source-to-sink path decomposition gives $|j_a^{(e)}|\le1$. Moreover
\[
 v_{\mathrm{tail}(e)}-v_{\mathrm{head}(e)}
 =\sum_a r_a(j_a^{(e)})^2>0,
\]
so $0\le j_e^{(e)}\le1$. Consequently
$|j_a^{(e)}-\one_{a=e}|\le1$ on every edge.

For a forcing $d_e$ only on edge $e$, the differentiated state equation
has $Ah=0$ and $\mathcal R h+d_e\one_e\in\operatorname{im}A^\top$.
Its unique solution is
\[
 h=\frac{d_e}{r_e}(j^{(e)}-\one_e).
\]
Indeed its divergence vanishes and its residual is
$(d_e/r_e)\mathcal R j^{(e)}$, a potential gradient. Uniqueness follows
by taking the scalar product of the homogeneous equation with its
circulation solution. Summing these solutions handles simultaneous
coefficient changes. The actual forcing on edge $e$ is
$\dot\beta_e x_e^\rho|x_e^\rho|$ in the symmetric case; in the
directional case it is $\dot\beta_e^+(x_e^\rho)^2$ for positive flow,
$-\dot\beta_e^-(x_e^\rho)^2$ for negative flow, and zero at zero.
Only one coefficient is active. In both cases
\[
 \frac{|d_e|}{r_e}\le
 \frac{q_e|x_e^\rho|^2}{2\beta_L|x_e^\rho|+\rho}
 \le\frac{q_eM}{2\beta_L}.
\]
Thus $\|\dot x^\rho(s)\|_\infty\le
M\sum_e q_e/(2\beta_L)$, and integration gives
\cref{eq:a-corr-sharper-flow} for the two regularized endpoint states.

Finally let $\rho\downarrow0$. At each endpoint the regularized flows
lie in the compact box $[-M,M]^E$. Every convergent subsequence has a
conserved limit. For any fixed conserved trial flow $y$, regularized
optimality gives
\[
 F_\beta(x^\rho)+\rho\|x^\rho\|_2^2/2
 \le F_\beta(y)+\rho\|y\|_2^2/2,
\]
where $F_\beta$ is the original primitive energy. Passing to the limit
shows the limit minimizes $F_\beta$; uniqueness makes it the original
physical flow. Hence the full endpoint sequences converge. Passing to
the limit in their uniform difference bound proves the assertion,
including zero and reversing flows. Uniform convergence over all $s$
is unnecessary for this endpoint argument.
\end{proof}

Resistance continuity is established: \citet[Theorem 3.12, pp.~34--35]{raber2022-thesis} proves
continuity of physical flows and normalized potentials in all positive
resistances by parametric convex optimization. The explicit linear
constant in \cref{eq:a-corr-linear-flow}, and its validity for separate
directional coefficients, provide the quantitative precision bound used
here. Qualitative continuity alone does not specify its bit budget.

\begin{corollary}[Strict margins and rational witness size]
\label{cor:a-corr-strict-margin}
For fixed rational nominations and independent rational positive
coefficient intervals (symmetric or directional), suppose a real profile
satisfies all rational signed arc bounds with common slack $\gamma>0$.
If $b\ne0$, rounding each coefficient within its own interval by at most
\[
 \delta\le\frac{\beta_L\gamma}{4mM}
\]
preserves at least $\gamma/2$ slack. For rational $\gamma$, a rational
feasible profile exists with encoding length polynomial in the input
length and $\max\{0,\log(1/\gamma)\}$. This is a conditional witness-size
and safe-rounding statement, not an algorithm for finding the initial
strictly feasible real profile.
\end{corollary}
\begin{proof}
Apply \cref{eq:a-corr-linear-flow}. For existence and encoding, round
coordinates to a dyadic grid of spacing at most $\delta$ and clip to the
rational interval endpoints. Clipping preserves the interval and cannot
increase distance from the original point in that interval. The number
of fractional bits and the integer parts are polynomially bounded as
claimed. The same construction applies separately to the two directional
coordinates. With $b=0$, every physical flow is zero and any rational
allowed profile is feasible whenever the capacities contain zero.
\end{proof}

The sharper estimate \cref{eq:a-corr-sharper-flow} preserves half
the slack whenever $\delta\le\beta_L\gamma/(mM)$. The radius in
\cref{cor:a-corr-strict-margin} also remains valid. The existence and size conclusion extends to any bounded rational
coefficient polytope $P$ within the same positive coefficient bounds,
including affine equalities, although arbitrary
coordinatewise rounding need not preserve its correlations. If $b=0$,
every physical flow is zero, and rational LP returns a polynomial-bit
profile in $P$ satisfying the same arc bounds. Otherwise choose
$\delta=2^{-k}$, where $k$ is the least nonnegative integer such that
\[
 2^{-k}\le\min\{1,\beta_L\gamma/(4mM)\}.
\]
Then $k$ is polynomially bounded in the input length and
$\max\{0,\log(1/\gamma)\}$. Choose a vector $a$ on the grid of spacing
$\delta/4$ within $\delta/4$ of the strictly feasible real profile
$\beta$, and form
\[
 P\cap\{z:\|z-a\|_\infty\le\delta/2\}.
\]
This is a nonempty bounded rational polytope because it contains
$\beta$. The chosen grid and the bounds on $\beta$ give $a$ and the box
inequalities polynomial encoding length in the original data and $k$.
A rational LP vertex therefore has polynomial bit length by determinant
bounds, and lies within
$3\delta/4$ of $\beta$. The sensitivity estimate preserves the required
physical slack. This proves existence of a short correlated witness;
if a sufficiently accurate center $a$ is supplied, LP recovers one.
It does not locate the unknown strictly feasible profile or its center.

For pressure bounds, let a fixed path between terminals have $k$ edges. On $[-M,M]$
every new directional law is $2\beta_U M$-Lipschitz, since its one-sided
derivatives have magnitude at most that number and the law is continuous
at zero. Thus splitting the coefficient and flow changes on each edge
and summing along the path gives
\begin{equation}
 \begin{split}
 |(\pi'_s-\pi'_t)-(\pi_s-\pi_t)|
 &\le k\bigl[\delta M^2+2\beta_U M\|x'-x\|_\infty\bigr]\\
 &\le kM^2\delta(1+4m\beta_U/\beta_L).
 \end{split}
 \label{eq:a-corr-pressure-rounding}
\end{equation}
Finitely many strict pressure-difference margins therefore survive
sufficiently fine rational rounding with the same polynomial-bit
principle. Absolute potential bounds require fixing a reference
potential, or explicitly choosing the additive gauge. With a fixed
reference, take paths from that vertex in the preceding estimate.

\subsection{Pressure filters and the conic boundary}

Even on a triangle, arbitrary upper pressure filters cannot be appended
as affine resistance restrictions to
\cref{thm:a-design-global-capacity}. Give its two-edge path total
resistance $A$ (each edge $A/2$) and its direct edge resistance $B$.
Under unit terminal transfer, equal drops and conservation yield
\[
 D(A,B)=\frac{AB}{(\sqrt A+\sqrt B)^2}.
\]
On the rational resistance segment joining $(A,B)=(1,9)$ and $(9,1)$,
the endpoint drops are $9/16$, while the midpoint $(5,5)$ has drop
$5/4$. Thus the upper restriction $D\le1$ accepts both endpoints and
rejects their midpoint, even though every physical flow lies in $[0,1]$.
The feasible resistance set is nonconvex and cannot be a projection of
a linear feasible set in these original coordinates.

For the opposite restriction under unit terminal transfer, $D\ge c$ is
a convex resistance condition by \cref{eq:a-corr-primitive}. It has an
exact extended conic representation: append $J\ge c/3$ to
\cref{eq:a-corr-cones,eq:a-corr-socp}. If such auxiliary variables exist,
weak duality gives $D\ge c$; conversely physical potentials and
\cref{eq:a-corr-physical-lift} attain $J=D/3$. Adding this inequality
may remove the strictly feasible conic center constructed above.
Consequently this representation does not by itself give an
unconditional exact rational feasible-output guarantee at a degenerate
boundary. It also concerns the single unit-transfer potential difference,
not arbitrary collections of potential bounds. These distinctions keep
the maximum-energy algorithm, the hard minimum-energy direction, and
the irrational capacity-witness example consistent.
